\documentclass[11pt]{article}
\usepackage[margin=1in]{geometry}
\usepackage{amsmath,amssymb}
\usepackage{enumitem}
\usepackage[T1]{fontenc}

\title{COMP2300 Set 8 Practice Exam}
\author{Topics: ILP, Dependences, Dynamic Scheduling, Renaming, ROB, Speculation, LSQ}
\date{}

\begin{document}
\maketitle

\noindent\textbf{Time:} 120 minutes \hfill
\textbf{Total:} 100 marks

\vspace{0.5em}
\noindent\textbf{Instructions}
\begin{itemize}[leftmargin=1.5em]
  \item Answer in English.
  \item Part A is multiple choice.
  \item Part B contains dependence analysis and out-of-order reasoning.
  \item Part C contains adapted past-exam style questions with source tags.
\end{itemize}

\section*{Part A: Multiple Choice (30 marks, 3 marks each)}
Choose exactly one option for each question.

\begin{enumerate}[label=\textbf{A\arabic*.}]
  \item Which dependence corresponds to a \texttt{RAW} hazard?
  \begin{enumerate}[label=(\Alph*)]
    \item Anti dependence
    \item Output dependence
    \item True dependence
    \item Control dependence
  \end{enumerate}

  \item Register renaming primarily eliminates:
  \begin{enumerate}[label=(\Alph*)]
    \item RAW only
    \item WAR and WAW
    \item Branch hazards only
    \item Cache misses
  \end{enumerate}

  \item An issue queue is used to:
  \begin{enumerate}[label=(\Alph*)]
    \item Store only committed results
    \item Hold waiting instructions until operands are ready
    \item Replace the branch predictor
    \item Encode immediate values
  \end{enumerate}

  \item Which statement about \texttt{ROB} is correct?
  \begin{enumerate}[label=(\Alph*)]
    \item It guarantees that instructions execute in order
    \item It supports in-order retirement of speculative results
    \item It stores branch target addresses only
    \item It removes all RAW hazards
  \end{enumerate}

  \item A precise exception requires that:
  \begin{enumerate}[label=(\Alph*)]
    \item Younger instructions may already be permanently committed
    \item Older instructions after the faulting instruction may commit first
    \item All older instructions have completed and the faulting instruction has not improperly committed
    \item The CPU flushes only the fetch stage
  \end{enumerate}

  \item Which dependence cannot be removed by renaming?
  \begin{enumerate}[label=(\Alph*)]
    \item WAR
    \item WAW
    \item RAW
    \item Output name dependence
  \end{enumerate}

  \item In a modern OOO design, speculative results are typically written first to:
  \begin{enumerate}[label=(\Alph*)]
    \item ARF only
    \item ROB
    \item Instruction memory
    \item BTB
  \end{enumerate}

  \item The main limitation of a scoreboard-only design highlighted in the lecture is:
  \begin{enumerate}[label=(\Alph*)]
    \item It has no ALUs
    \item It cannot support any form of data dependence
    \item It struggles with false dependences and precise recovery
    \item It eliminates branch hazards completely
  \end{enumerate}

  \item A load-store queue is used because:
  \begin{enumerate}[label=(\Alph*)]
    \item Loads and stores are always safe to reorder
    \item Memory addresses may alias, so load/store reordering needs checking
    \item It replaces all caches
    \item It is the same thing as a branch history table
  \end{enumerate}

  \item Memory-level parallelism refers to:
  \begin{enumerate}[label=(\Alph*)]
    \item Using one memory request at a time very quickly
    \item Exploiting multiple memory requests in flight concurrently
    \item Replacing DRAM with SRAM
    \item Turning branches into arithmetic
  \end{enumerate}
\end{enumerate}

\section*{Part B: Computational / Reasoning Problems (50 marks)}
\begin{enumerate}[label=\textbf{B\arabic*.}]
  \item \textbf{Classifying dependences (8 marks)}\\
  For each pair below, classify the dependence as \textbf{RAW}, \textbf{WAR}, \textbf{WAW}, or \textbf{none}.
\begin{verbatim}
i1: LDR R2, [R6, #0]
i2: ADD R3, R2, R5
\end{verbatim}
\begin{verbatim}
i3: ADD R6, R3, R5
i4: LDR R2, [R6, #0]
\end{verbatim}
\begin{verbatim}
i5: ADD R2, R7, R8
i6: SUB R2, R9, R10
\end{verbatim}

  \item \textbf{Renaming to remove false dependences (8 marks)}\\
  Rename the following code using temporary names \texttt{T0} and \texttt{T1} so that false dependences are removed while true dependences are preserved.
\begin{verbatim}
i1: ADD R1, R2, R3
i2: SUB R2, R4, R5
i3: ADD R6, R1, R7
i4: MUL R2, R8, R9
\end{verbatim}
  State which original false dependence(s) you removed.

  \item \textbf{ROB allocation and retirement order (8 marks)}\\
  Suppose three instructions enter rename in program order and are assigned tags \texttt{ROB0}, \texttt{ROB1}, and \texttt{ROB2}.
  \begin{enumerate}[label=(\alph*)]
    \item Which tag belongs to the oldest instruction?
    \item If \texttt{ROB2} finishes execution first, may it retire first? Why or why not?
    \item Which instruction must be at the head of the ROB before \texttt{ROB2} may retire?
  \end{enumerate}

  \item \textbf{Speculation and precise exceptions (8 marks)}\\
  An older load instruction causes an exception. Two younger arithmetic instructions have already finished execution out of order.
  \begin{enumerate}[label=(\alph*)]
    \item In a design that writes results directly into the architectural register file, why can this become an imprecise exception?
    \item How does a ROB-based design avoid this problem?
  \end{enumerate}

  \item \textbf{Load/store reordering (8 marks)}\\
  For each case below, state whether reordering is safe and briefly justify:
  \begin{enumerate}[label=(\alph*)]
    \item An older store writes address \texttt{A}; a younger load reads address \texttt{B}, where \texttt{A} and \texttt{B} are known to be different.
    \item An older store writes address \texttt{A}; a younger load reads address \texttt{A}.
    \item The address of the older store is not yet known.
  \end{enumerate}

  \item \textbf{OOO overlap on a cache miss (10 marks)}\\
  A load instruction misses in cache and waits many cycles for memory. Behind it are:
  \begin{itemize}[leftmargin=1.5em]
    \item one dependent ADD
    \item two independent arithmetic instructions
    \item one younger branch with operands already ready
  \end{itemize}
  \begin{enumerate}[label=(\alph*)]
    \item Which instruction must wait because of a true dependence?
    \item Which younger instructions may still issue in an OOO design?
    \item Why is this impossible in a strict in-order design?
  \end{enumerate}
\end{enumerate}

\section*{Part C: Past Exam Questions (Source-Tagged) (20 marks)}
\begin{enumerate}[label=\textbf{C\arabic*.}]
  \item \textbf{Register renaming (Source: exam24r.pdf, Q7b; 2023\_Final\_Exam.pdf, Q2 rename topic, adapted) (7 marks)}\\
  Explain how register renaming removes false dependences while preserving true dependences. Why does this help out-of-order execution?

  \item \textbf{OOO pipeline issues addressed by ARF+ROB (Source: exam2025-main.pdf, Q3C; 2023\_Final\_Exam\_SuppDA.pdf, ARF+ROB discussion, adapted) (7 marks)}\\
  List four problems in a scoreboard-style out-of-order pipeline that are addressed by adding register renaming and hardware speculation.

  \item \textbf{Issue queue and dynamic scheduling (Source: exam2025-main.pdf, Q3; exam24r.pdf, Q7c, adapted) (6 marks)}\\
  Briefly explain the role of the issue queue in a dynamically scheduled processor. Why does it help tolerate long-latency misses better than a simple in-order pipeline?
\end{enumerate}

\end{document}
